\scanpage{208}
\[
G_k,H_k\subset\Omega_k
\]
\[
G_k\cap H_k=\{u_\lambda\mid\lambda\in S_k,\ \lambda T=T\ \text{i.e. }(\lambda\{T\}-1)T=0\}
\]
Si on pose $\lambda\{T\}=\sum_{i\geq0}\lambda_i T^{(i)}$, alors la condition sur $\lambda$ envisagée est
\[
\boxed{\lambda_0=1}\quad\deleted{\lambda_0=1},\qquad\boxed{(i+1)\lambda_i=0\qquad\forall i\geq1}
\]
\deleted{Ce qui implique en particulier $\lambda_0=1$.}
\deleted{\unsure{$k$ de car. $0$, cela implique plus que}}
Si $k$ sans torsion, cela implique $\lambda_i=0$ pour $i\geq1$ i.e. $\lambda=1$, i.e. $u_\lambda=\operatorname{id}$, d'où $G_k\cap H_k\inserted{=\unsure{(e)|\{e\}}}$ dans ce cas, --- et dans ce cas seulement.
[Dans ce cas, \inserted{$k$ de car. $0$ i.e. une $\unsure{Q|\mathbb Q}$-alg.,} $G_k$ est aussi le groupe \inserted{$G'_k$} de tous les endom. de $S_k=k\{T\}$ comme $k$-alg. \uline{augmentée} \deleted{\unsure{p.d.}} et par suite \deleted{l'hom. \unsure{dans} $\Omega$}
\diagram{\begin{tikzcd}[ampersand replacement=\&, column sep=small]
(\unsure{1|0}\arrow[r] \& H'_k\arrow[r] \& )\,\Omega_k\arrow[r,dash] \& G'_k=\operatorname{End}_{\substack{k\text{-alg}\\\deleted{\mathrm{aug}}}}(\underset{S_k}{\underbrace{k\{T\}}})\\
\&\&\&G_k\arrow[u,"(??)"]\arrow[ul,bend left=28]\arrow[ul,bend right=18,dash]
\end{tikzcd}}
donne $\Omega_k$ comme produit $\frac12$ direct de $G_k\inserted{\simeq G'_k}$ et de $H'_k$\unsure{\ldots} ]

\deleted{Dans le cas général, on peut dire que}
\diagram{\begin{tikzcd}[ampersand replacement=\&, column sep=small]
1\arrow[r,dash] \& H_k\arrow[r,dash] \& \unsure{\Omega_k}\arrow[r] \& \unsure{G'_k|G_k}=\operatorname{End}_{k\text{-alg. augm.}}(k\{T\})\\
\&\&\&G_k\arrow[u]\arrow[ul,bend left=28,dash]\arrow[ul,bend right=18,dash]\\
\&\&\&G_k\cap H_k\simeq\{u_\lambda\mid\lambda=1+\sum_{i\geq1}\lambda_iT^{(i)},\ (i+1)\lambda_i=0\}\arrow[u]
\end{tikzcd}}
\margin{207}
